\section{Experimental Setup}

We design experiments to test whether uncertainty probes generalize across LLM families. Specifically, we address three research questions:

\textbf{RQ1:} Does a probe trained on one model family detect hallucinations when applied to another family's hidden states?

\textbf{RQ2:} Does affine alignment enable transfer between models with different hidden dimensions?

\textbf{RQ3:} What is the maximum transfer gap across all family pairs, and does it stay below the 0.10 threshold?

\subsection{Dataset}

\textbf{TruthfulQA} \citep{lin2022truthfulqa}. We evaluate on the generation subset containing 817 questions designed to elicit plausible but incorrect answers from language models. The benchmark covers 38 categories including health, law, finance, and politics. Ground truth correctness labels enable direct AUROC computation for hallucination detection.

We split the dataset 80/20 into training (653 questions) and validation (164 questions) using a fixed seed (42). The training split is used to compute semantic entropy labels and train probes; the validation split is held out for all transfer evaluations.

\begin{table}[h]
\centering
\caption{Dataset splits and usage.}
\label{tab:splits}
\begin{tabular}{@{}lrl@{}}
\toprule
Split & Questions & Usage \\
\midrule
Train & 653 & SE label computation, probe training, aligner fitting \\
Validation & 164 & Transfer evaluation (all reported metrics) \\
\bottomrule
\end{tabular}
\end{table}

\subsection{Models}

We test three instruction-tuned LLMs from different organizations:

\begin{table}[h]
\centering
\caption{Model specifications.}
\label{tab:models}
\begin{tabular}{@{}lllrr@{}}
\toprule
Model & Provider & Params & Hidden Dim & Layers \\
\midrule
Llama-3-8B-Instruct & Meta & 8B & 4096 & 32 \\
Mistral-7B-Instruct-v0.2 & Mistral AI & 7B & 4096 & 32 \\
Qwen-2-7B-Instruct & Alibaba & 7B & 3584 & 28 \\
\bottomrule
\end{tabular}
\end{table}

\subsection{Baselines}

\textbf{Multi-sample Semantic Entropy} \citep{farquhar2024detecting}. The gold-standard for hallucination detection. For each question, we generate 5 responses at temperature 0.7, cluster them by semantic equivalence using DeBERTa-v3-large-mnli-fever-anli-ling-wanli NLI model, and compute entropy over cluster distributions.

\textbf{Token-level Entropy.} Average entropy of next-token distributions during generation. This cheap baseline achieves AUROC $\sim$0.55 and serves as a lower bound.

\textbf{Direct (Untransferred) Probe.} Each model's SEP evaluated on its own hidden states. This provides the baseline AUROC from which we measure transfer gaps.

\subsection{Evaluation Metrics}

\textbf{AUROC.} Primary metric for hallucination detection. Measures the probability that a randomly chosen correct answer has lower predicted uncertainty than a randomly chosen incorrect answer.

\textbf{Transfer Gap.} For source model $i$ and target model $j$:
\begin{equation}
\text{gap}_{i \to j} = \text{AUROC}_{i \to i} - \text{AUROC}_{i \to j}
\end{equation}

\textbf{Success Criteria.} Pre-registered thresholds:
\begin{itemize}
\item Mean gap $< 0.05$: Strong evidence for architecture-invariant encoding
\item Max gap $< 0.10$: No pair shows prohibitive transfer degradation
\end{itemize}
